\section{Winding metadata and fault-current phasors under a documented switch outage}
\label{sec:p16-winding-fault}
Using the pinned OpenDSS IEEE 123-bus feeder revision \texttt{5005c668a72d20775f4c2d060feebb2866ba1d38} and \texttt{dss-python==0.15.7}, we extended the named-switch sensitivity to extract transformer winding connection codes, nominal winding voltages, apparent-power ratings, tap settings, terminal bus labels, and complex element-terminal currents. Eight transformer element records were found. The extraction does not, by itself, establish a complete phase/node transfer matrix, grounding impedances, or source-reachable energized phases; delta/wye coupling and regulator winding behavior cannot be replaced by one-to-one graph edges.

Four phase-1 single-line-to-ground fault probes at resistance \(0.1\,\Omega\) were solved at buses 60 and 83 for the published feeder and for a fresh compilation with \texttt{Line.sw1} disabled. All four numerical solves converged. Under the published state, the fault-bus phase-1 voltage magnitudes were \(0.13718528\) pu at bus 60 and \(0.07545517\) pu at bus 83. Under the disabled-switch state both were reported as zero. The latter state is therefore not an independently demonstrated energized protection scenario. In particular, fault-current comparisons between an energized case and a zero-voltage case cannot establish a robust directional relay margin.

Complex terminal-current arrays were retained for candidate elements \texttt{Line.l58}, \texttt{Line.l55}, \texttt{Line.l84} and \texttt{Line.l82}, using the OpenDSS \emph{into-element} current convention. No inference is made about installed CT orientation, polarity, primary/backup zone boundaries or field relay settings. Nominal winding kVA metadata do not alone certify line ampacity or dynamic thermal limits. Likewise, converged fault studies do not verify inverter electromagnetic transients or switching permissions.

The pinned CI workflow completed twelve adversarial regression tests and an independent structural audit of the four model scenarios. These results are a reproducible \emph{model sensitivity}, not a physical restoration authorization. The open questions are phase-specific source reachability through transformer/regulator windings, grounding transfer, thermal adequacy and genuine pre/post restoration cases in which both states remain energized.